%% Complete elementary proofs for the active CPL substitution/replacement pair.
\begin{annexitem}{annex:p2.prop.substitution}{thm:p02-uniform-substitution}{Semantic uniform substitution}
\begin{proof}
Fix \(\upsilon\) and a valuation \(v\).  Each \(\upsilon(p_n)\) is a formula,
so its value under \(v\) is exactly one of \(\mathbf F,\mathbf T\).
Consequently the formula in
Definition~\ref{def:p02-uniform-substitution} defines a valuation
\(v_\upsilon\) on the entire variable set.

We prove
\[
 \widehat v\bigl(\varphi[\upsilon]\bigr)
    =\widehat{v_\upsilon}(\varphi)
\]
by structural induction on the original formula \(\varphi\).
All primitive formation cases are as follows.

For a variable \(p_n\), the left side is
\(\widehat v(\upsilon(p_n))=v_\upsilon(p_n)\), which is the right side.
For \(\bottom\), substitution leaves the formula unchanged and both sides
are \(\mathbf F\).

For \(\varphi=\neg\alpha\), the substituted formula is
\(\neg(\alpha[\upsilon])\).  Its value is \(\mathbf T\) exactly when
\(\widehat v(\alpha[\upsilon])=\mathbf F\).
By the induction hypothesis this holds exactly when
\(\widehat{v_\upsilon}(\alpha)=\mathbf F\), which is the condition for
\(\widehat{v_\upsilon}(\neg\alpha)=\mathbf T\).
Since there are exactly two truth values, the two negation values agree.

For \(\varphi=(\alpha\land\beta)\), its substituted formula has value
\(\mathbf T\) exactly when both
\(\widehat v(\alpha[\upsilon])=\mathbf T\) and
\(\widehat v(\beta[\upsilon])=\mathbf T\).
The two induction hypotheses turn these into
\(\widehat{v_\upsilon}(\alpha)=\mathbf T\) and
\(\widehat{v_\upsilon}(\beta)=\mathbf T\), respectively.
Their conjunction is exactly the condition for
\(\widehat{v_\upsilon}(\alpha\land\beta)=\mathbf T\); hence the values agree.

For \(\varphi=(\alpha\lor\beta)\), its substituted formula has value
\(\mathbf T\) exactly when at least one of
\(\widehat v(\alpha[\upsilon])=\mathbf T\) and
\(\widehat v(\beta[\upsilon])=\mathbf T\) holds.
By the induction hypotheses this is exactly the condition that at least one of
\(\widehat{v_\upsilon}(\alpha)=\mathbf T\) and
\(\widehat{v_\upsilon}(\beta)=\mathbf T\) holds.
That condition defines
\(\widehat{v_\upsilon}(\alpha\lor\beta)=\mathbf T\), so the values agree.

For \(\varphi=(\alpha\to\beta)\), its substituted formula has value
\(\mathbf F\) exactly when
\(\widehat v(\alpha[\upsilon])=\mathbf T\) and
\(\widehat v(\beta[\upsilon])=\mathbf F\).
The induction hypotheses give precisely
\(\widehat{v_\upsilon}(\alpha)=\mathbf T\) and
\(\widehat{v_\upsilon}(\beta)=\mathbf F\).
By the implication clause, their conjunction is exactly the condition for
\(\widehat{v_\upsilon}(\alpha\to\beta)=\mathbf F\).
Bivalence again gives equality of the two values.  This exhausts the primitive
grammar and proves the identity.

If \(\entails\varphi\), then every valuation, including \(v_\upsilon\),
assigns \(\mathbf T\) to \(\varphi\).  The identity gives
\(\widehat v(\varphi[\upsilon])=\mathbf T\).  Since \(v\) was arbitrary,
\(\entails\varphi[\upsilon]\).

If \(\varphi\equiv\psi\), then
\(\widehat{v_\upsilon}(\varphi)=\widehat{v_\upsilon}(\psi)\).
Applying the identity to each side gives
\[
 \widehat v(\varphi[\upsilon])
 =\widehat{v_\upsilon}(\varphi)
 =\widehat{v_\upsilon}(\psi)
 =\widehat v(\psi[\upsilon]).
\]
This holds for every \(v\), proving
\(\varphi[\upsilon]\equiv\psi[\upsilon]\).
Neither deduction uses a converse or surjectivity of
\(v\mapsto v_\upsilon\).
\end{proof}
\end{annexitem}

\begin{annexitem}{annex:p2.prop.replacement}{thm:p02-replacement}{Equivalence congruence and selected-occurrence replacement}
\begin{proof}
Reflexivity follows because
\(\widehat v(\alpha)=\widehat v(\alpha)\) for every \(v\).
If \(\alpha\equiv\beta\), symmetry of equality gives
\(\widehat v(\beta)=\widehat v(\alpha)\) for every \(v\), so
\(\beta\equiv\alpha\).  If also \(\beta\equiv\gamma\), transitivity of
equality gives \(\widehat v(\alpha)=\widehat v(\gamma)\) for every \(v\),
so \(\alpha\equiv\gamma\).

For the nullary connective, every valuation assigns \(\mathbf F\) to
\(\bottom\), giving \(\bottom\equiv\bottom\).
Assume \(\alpha\equiv\alpha'\), and, for the binary cases, also
\(\beta\equiv\beta'\).  Fix any valuation \(v\).

The value of \(\neg\alpha\) is \(\mathbf T\) exactly when
\(\widehat v(\alpha)=\mathbf F\).  Equality of the input values makes this
equivalent to \(\widehat v(\alpha')=\mathbf F\), the condition for
\(\widehat v(\neg\alpha')=\mathbf T\).  Thus
\(\neg\alpha\equiv\neg\alpha'\).

The value of \(\alpha\land\beta\) is \(\mathbf T\) exactly when both
\(\widehat v(\alpha)=\mathbf T\) and \(\widehat v(\beta)=\mathbf T\).
Input equality preserves these two conditions, so it preserves their
conjunction and the truth value of the conjunction formula.

The value of \(\alpha\lor\beta\) is \(\mathbf T\) exactly when at least one of
\(\widehat v(\alpha)=\mathbf T\) and \(\widehat v(\beta)=\mathbf T\) holds.
Input equality preserves each condition and hence their disjunction.
Thus the two disjunction formulas have equal truth values.

The value of \(\alpha\to\beta\) is \(\mathbf F\) exactly when
\(\widehat v(\alpha)=\mathbf T\) and \(\widehat v(\beta)=\mathbf F\).
Input equality preserves both conditions and hence their conjunction.
Thus the two implication formulas have equal truth values.
Bivalence justifies equality in each case from agreement on the condition
for one of the two values.  Since \(v\) was arbitrary, these arguments prove
the three binary congruences.  All primitive connectives have been checked.
The paragraphs up to this point prove both parts of
Exercise~\ref{ex:prop-metatheory-09}, including the nullary case
\(\bottom\equiv\bottom\);
Convention~\ref{def:p02-w02-p02-004-exercise-correspondence} matches them
with the steps of the feedback sketch for that exercise.

We now justify the occurrence operation and its preservation claim.
Let \(c(\alpha)\) count occurrences of \(\neg,\land,\lor,\to\), with
variables and \(\bottom\) assigned complexity zero as in
Section~\ref{sec:prop-syntax-semantics-01}.
An immediate subformula has strictly smaller complexity:
\(c(\alpha)<c(\neg\alpha)=c(\alpha)+1\), and both
\(c(\alpha)\) and \(c(\beta)\) are smaller than
\(c(\alpha\circ\beta)=c(\alpha)+c(\beta)+1\).
Consequently, if one position properly extends another, its subformula has
strictly smaller complexity.  Two positions whose subformulas are both the
same string \(\psi\) therefore cannot properly contain one another.
In particular, if \(\epsilon\in O\), then \(\varphi=\psi\) and
\(O=\{\epsilon\}\).  If the root is unselected, the sets \(O_i\) select exactly
the corresponding occurrences inside each immediate subformula.

The replacement recursion is thus well-defined on the original finite parse
tree.  A selected root returns the formula \(\chi\) without further
recursion.  An unselected leaf returns that leaf; an unselected negation or
binary node combines recursively obtained formulas by its original primitive
constructor.  Every recursive call is on a proper subformula, so recursion
terminates and yields a formula of the same language.

Assume \(\psi\equiv\chi\).  We prove
\(\varphi\equiv\varphi[O\leftarrow\chi]\) by structural induction on
\(\varphi\), allowing every admissible chosen set \(O\) at each step.
If its root is selected, the result is \(\chi\) and
\(\varphi=\psi\), so the claim is the hypothesis \(\psi\equiv\chi\).

Suppose the root is unselected.  For a variable or \(\bottom\), its only
position is the root, so \(O=\varnothing\) and the output is unchanged.
Reflexivity proves these two cases.

For \(\varphi=\neg\alpha\), the induction hypothesis gives
\(\alpha\equiv\alpha[O_1\leftarrow\chi]\).
Negation congruence therefore gives
\(\neg\alpha\equiv\neg(\alpha[O_1\leftarrow\chi])\), exactly the required
replacement result.

For \(\varphi=(\alpha\land\beta)\), the induction hypotheses give equivalence
of \(\alpha\) and its \(O_1\)-replacement and of \(\beta\) and its
\(O_2\)-replacement.  Conjunction congruence proves the result.
For \(\varphi=(\alpha\lor\beta)\), the same two induction hypotheses together
with the separately proved disjunction congruence prove the result.
For \(\varphi=(\alpha\to\beta)\), those induction hypotheses and implication
congruence prove the result.  These are all remaining primitive constructors.

The empty choice set is included: at every node it follows the unselected
case and leaves the whole original formula unchanged.  Selection of all
occurrences is also included, since all selected positions belong to the
original finite tree and no inserted formula is revisited.
\end{proof}
\end{annexitem}
